\section{Controlled Output Interventions}\label{app:causal}
The five interventions per prospective window were specified after the primary results were inspected and before intervention scoring. They hold weights, thresholds, proposals, and original predictions fixed. Their effects are conditional on those decisions.

\paragraph{Historical-label controls.}
``Ours new'' copies Ours appended sources, coordinates, labels, and ranking metadata. Replacing only the historical-label prefix with either baseline or unconditional labels measures the effect of old assignments at fixed new geometry. Old-point coordinates and ordering are unchanged.

\paragraph{New-point controls.}
``Witness'' keeps directly mask-supported parent additions with their parent labels, paired with either preserved or authorized history. ``Nearest'' fixes Ours history and assigns supported additions to the unique nearest seed identity without the second seed-frame requirement on new points. Historical authorization, including its local check, remains fixed. All controls share the Ours classification metadata and postprocessing.

\begin{center}\small
\begin{tabular}{@{}lrrr@{}}
\toprule
Export intervention & 0009 & 0010 & Mean \\
\midrule
Unconditional old IDs + Ours new & 23.253 & 11.420 & 17.337 \\
Preserved old IDs + Ours new & 23.166 & 14.198 & 18.682 \\
Preserved old IDs + witness & 23.089 & 12.809 & 17.949 \\
Authorized old IDs + witness & 23.230 & 12.809 & 18.019 \\
Authorized old IDs + nearest & 23.327 & 14.198 & 18.762 \\
\textbf{Ours (fixed)} & 23.327 & 14.198 & 18.762 \\
\bottomrule
\end{tabular}

\captionof{table}{All controlled output interventions, AP50:90 in percent. Only the operations specified above vary.}
\end{center}

For 0009, the nearest-seed control retains 35 sources rejected by Ours and differs at five reference positions, yet AP is equal. For 0010, its source set, labels, and reference predictions are identical to Ours. The witness-history versus witness-preservation contrast is +0.141 / +0.000 AP points in 0009/0010. These paired effects depend on which additions are held fixed and need not sum to the total baseline-to-Ours change.

\subsection{Interpreting historical edits}
Each baseline predicted ID is matched to the same-class canonical reference instance with maximum confidence-weighted intersection; ties use the lower reference ID. Multiple fragments may match the same instance, and IDs without same-class overlap remain unmapped. This correspondence is then held fixed across variants. Each reference point retains its nearest baseline source within 0.20\,m, and only the historical label at that source changes.

For mapped endpoints, wrong-to-correct transitions are corrections, correct-to-wrong transitions are damage, and correct-to-correct transitions consolidate fragments of the same reference object. Wrong-to-wrong transitions remain unresolved. Missing endpoint correspondences and void reference points are reported separately. Counts other than ``Old edits'' are confidence-weighted reference mass, rather than source-point accuracy.

\begin{center}\footnotesize
\begin{tabular}{@{}llrrrrrr@{}}
\toprule
Window & Method & Old edits & Correction & Damage & Consolidation & Unresolved & Unknown \\
\midrule
0009 & Unconditional & 26,788 & 101.1 & 1072.9 & 10848.0 & 7152.8 & 3348.2 \\
0009 & \textbf{Ours} & 12,855 & 0.0 & 0.0 & 10671.7 & 2.7 & 0.0 \\
0010 & Unconditional & 3,427 & 0.0 & 0.0 & 0.0 & 0.0 & 3552.2 \\
0010 & \textbf{Ours} & 0 & 0.0 & 0.0 & 0.0 & 0.0 & 0.0 \\
\bottomrule
\end{tabular}

\captionof{table}{Historical outcomes under a fixed baseline-to-reference correspondence. Unknown endpoints are excluded from correct/incorrect categories.}
\end{center}

The authorized 0009 transitions are $27\!\to\!17$ (8,745 sources, building 2168) and $36\!\to\!48$ (4,110 sources, building 2486). Both consolidate fragments of the same object; 2.736 reference mass remains wrongly assigned. The unconditional 0010 transition targets ID 6, which is absent from the baseline correspondence, leaving 3,552.226 reference mass unknown in this diagnostic. Its effect on AP is measured by the fixed-new-geometry intervention in Table~\ref{tab:causal}.

\section{Runtime and External Comparison}\label{app:runtime}
\subsection{Revision runtime}
On an AMD Ryzen 9 9950X, cached graph assignment takes 117.66\,ms for 0009 and 17.33\,ms for 0010. The graph timer includes both native multi-source distance searches and the certificate calculation. The pair timer evaluates merge decisions from cached per-view evidence. Model inference, projection, proposal construction, loading, and export are excluded. These are single CPU runs rather than end-to-end latency estimates.

\begin{center}\small
\begin{tabular}{@{}lrrrrrr@{}}
\toprule
Window & Nodes & Edges & Seeds & Graph (ms) & Pairs (ms) & Prior prep. (s) \\
\midrule
0009 & 42,955 & 457,332 & 61,886 & 117.66 & 0.180 & 5.201 \\
0010 & 15,091 & 129,426 & 11,665 & 17.33 & 0.001 & 2.323 \\
\bottomrule
\end{tabular}

\captionof{table}{Revision input sizes and component timings. Edges are directed. Prior preparation was measured in a separate run and is not added to the replay timers.}
\end{center}

The replay driver's peak resident memory is 198.4\,MiB across the two cached runs. Map export takes 0.286--0.290\,s for 0009 and 0.088--0.089\,s for 0010. Native distances and certificate arrays reproduce the archived results exactly. The geometry weights were fitted on exposed 0000/0002 data; neither the replay nor the extension refits them.

\subsection{MaskClustering core adaptation}\label{app:external}
We run the official MaskClustering graph construction, iterative view-consensus clustering, and postprocessing code through a CPU adapter. Six upstream files are pinned; only device substitutions and an unused backprojection import are changed. Inputs are disjoint highest-score SAM3 query supports on the same measured sources, using the causal visibility cache. Query scores provide the building/car readout. This replaces the native CropFormer~\cite{cropformer}/OpenCLIP~\cite{openclip} frontend and also differs from the internal backbone's semantic readout.

We retain the upstream ScanNet defaults: visible and undersegmentation ratios 0.3, consensus 0.9, contained-mask ratio 0.8, point-filter ratio 0.5, DBSCAN radius 0.1\,m with four minimum points, two supporting masks, and overlap merge ratio 0.8. Empty graphs export no clusters. Source ties use stable object order in the nearest-map readout.

\begin{center}\small
\begin{tabular}{@{}lrrr@{}}
\toprule
Window & Primary AP & Canonical AP & Uniform-score AP \\
\midrule
0009 & 0.0071 & 0.0082 & 0.1664 \\
0010 & 0.0000 & 0.0000 & 0.0000 \\
\bottomrule
\end{tabular}

\captionof{table}{MaskClustering core adaptation on the prospective windows. AP is reported in percent; the four-window extension is included in Table~\ref{tab:results}b.}
\end{center}

The adapter yields 2,542 / 662 intermediate clusters and 98 / 23 final objects on 0009/0010. Fragmentation is severe and AP is low. Because the inputs and readout have changed and indoor defaults are retained, this comparison diagnoses the adaptation; it cannot establish superiority over the native MaskClustering pipeline.
